%% Dans ce fichier, les �ventuels accents sont cod�s
%% en ISO-8859-15 (aka latin9)
%% La ligne suivante donne trois exemples 
%% e accent aigu : �, e accent grave : �, e-dans-l'o : �
%%%%%ENTREES%%%%
\begin{tabarg}
\multicolumn{3}{l}{\titretab {Entr\'ees}}\\
\hline
$X$ & Vecteur, \texttt{float} & S�rie d'origine (unit� indiff�rente)\\
$Xflag$ & Vecteur, \texttt{int} & vecteur de flag associ� � X\\
$flag_{inf}$ & Coeff., \texttt{float} & Borne inf�rieure de la plage des valeurs de flag � remplacer.\\
$flag_{sup}$ & Coeff., \texttt{float} & Borne sup�rieure de la plage des valeurs de flag � remplacer.\\
$val_{remplac}$  &  Coeff, \texttt{float} & Valeur de remplacement \\
\end{tabarg}

\medskip
%%%%%SORTIES%%%%
\begin{tabarg}
\multicolumn{3}{l}{\titretab {Sorties}}\\
\hline
$X_{mod}$  & Vecteur, \texttt{float} & S�rie des donn�es modifi�es  (m�me unit� que $X$)\\
\end{tabarg}
